\newglossaryentry{lehen}{
    type=terms,
    name={\emph{lêhen}},
    description={Mhd.: \enquote{Lehen}, \enquote{Leihe}. Das Lehen bezeichnet im Feudalsystem die Übertragung von Land, Amt oder Rechten durch einen Lehnsherrn an einen Vasallen gegen Eid, Treue und Dienst. Das Lehenswesen bildete das Rückgrat der mittelalterlichen Gesellschaftsordnung; fiel ein Lehnsträger ohne Erben, kehrte das Lehen an den Herrn zurück}
}

\newglossaryentry{Raabs an der Thaya}{
    type=terms,
    name={Raabs an der Thaya},
    description={Kleinstadt im Waldviertel des Herzogtums \emph{Oesterr\^{\i}che}, an der Grenze zu M\"ahren und B\"ohmen. Siehe \nameref{raabs}}
}

\newglossaryentry{Oesterriche}{
    type=terms,
    name={\emph{Oesterr\^{\i}che}},
    description={Mhd. Bezeichnung f\"ur das Herzogtum \enquote{\"Osterreich}. Im 13.~Jahrhundert war dies die gebr\"auchliche Raumbezeichnung; der moderne Verwaltungsbegriff Nieder\"osterreich existierte in dieser Form noch nicht}
}

\newglossaryentry{Zwettl}{
    type=terms,
    name={Zwettl},
    description={Stadt im Waldviertel und Sitz des 1138 gegr\"undeten Zisterzienserstifts Zwettl, das im 13.~Jahrhundert als Grundherrschaft und geistliches Zentrum von hoher regionaler Bedeutung war}
}

\newglossaryentry{Baden}{
    type=terms,
    name={Baden},
    description={Baden bei Wien, Markt- und Thermalort im Herzogtum \emph{Oesterr\^{\i}che}; durch seine Lage am Ostrand des Wienerwaldes und nahe wichtigen Verkehrswegen im Hochmittelalter regional bedeutsam}
}

\newglossaryentry{Rauheneck}{
    type=terms,
    name={Rauheneck},
    description={Burg bei Baden im Wienerwald, im 12. und 13.~Jahrhundert Sitz eines Ministerialengeschlechts, das im Umfeld landesf\"urstlicher Herrschaft in \emph{Oesterr\^{\i}che} belegt ist}
}

\newglossaryentry{ordo mundi}{
    type=terms,
    name={\emph{ordo mundi}},
    description={Lat.: \enquote{Ordnung der Welt}. Die mittelalterliche Vorstellung einer von Gott eingesetzten, unveränderlichen Weltordnung, die alle Bereiche des Lebens — Natur, Gesellschaft, Kirche — durchdringt und in ein hierarchisches Gefüge einbettet. Abweichungen von dieser Ordnung galten als Sünde und zogen Unheil nach sich}
}

\newglossaryentry{liberum arbitrium}{
    type=terms,
    name={\emph{liberum arbitrium}},
    description={Lat.: \enquote{freier Wille}. Zentraler Begriff der mittelalterlichen Theologie und Philosophie (u.\,a. Augustinus, Thomas von Aquin). Der freie Wille des Menschen gilt als Gabe Gottes, durch die der Mensch zwischen Gut und Böse, Tugend und Laster zu wählen vermag; er bildet die Grundlage moralischer Verantwortung}
}

\newglossaryentry{triuwe}{
    type=terms,
    name={\emph{triuwe}},
    description={Mhd.: \enquote{Treue}, \enquote{Loyalität}, \enquote{Zuverlässigkeit}; umfasst die Gesamtheit der sozialen und moralischen Bindungen, die ein Mensch gegenüber Herren, Gefährten und der Gemeinschaft eingeht. Im ritterlichen Tugendsystem einer der höchsten Werte}
}

\newglossaryentry{Interregnum}{
    type=terms,
    name={\emph{Interregnum}},
    description={Lat.: \enquote{Zwischenherrschaft}. Historische Bezeichnung für die Zeit ohne anerkannten deutschen König/Kaiser (1254--1273), in der Reichsgewalt und Rechtsfrieden weitgehend verfielen. In der Wahrnehmung der Zeitgenossen eine Epoche des Chaos, der Fehden und der Rechtlosigkeit}
}
